% !TeX root = chapter_3_main.tex
\section{Molasses temperature optimization}

Before performing any optimization, the six beams in the vacuum chamber were realigned in order to minimize their overlap mismatch at the MOT position. This procedure was necessary since thermal drifts can induce misalignment in the optical components.

Preliminary values for the compensation coil currents has also been obtained prior the realignment using the Release-and -Recapture (R\&R) technique:
\begin{equation}
    \begin{aligned}
        I_{\mathrm{x}} &= 80~\mathrm{mA} \\
        I_{\mathrm{y}} &= -200~\mathrm{mA} \\
        I_{\mathrm{z}} &= -400~\mathrm{mA}.
    \end{aligned}
    \label{eq:shim_preliminary_values}
\end{equation}
These values were used as the starting point for the subsequent optimization procedure.
%! Ricorda di aggiungere alla sezione dei metodi la considerazione sul t_flash


%*###############################################################
\subsection{Detuning frequency optimization}
%*###############################################################
The first optimization step consisted of measuring the molasses temperature for different values of the cooling laser detuning. The detuning was varied from $-24.4$~MHz (4.0~$\Gamma/2\pi$) to $-184$~MHz (30.4~$\Gamma/2\pi$), with a step of $-11.4$~MHz.


The resulting average temperature is shown in Figure \ref{fig:detuning_molasses_temperature} as a function of the detuning expressed in units of $\Gamma/2\pi$. At small detunings, the data follow the expected $1/|\delta|$ dependence typical of the Sisyphus cooling. As the detuning increases two phenomena becomes relevant: excessive decrease of the optical cycle rate $\Gamma_p \propto 1/\delta^2$ and the approach of the transition $F=2 \rightarrow F'=2$ which is found $267$~MHz below the cooler target transition $F=2 \rightarrow F'=2$ {\color{blue}(see Figure ref??)}. These effects cause the temperature curve to flatten in the region $|\delta|\approx15-20$ and eventually start to increase again for larger detunings.

\begin{figure}[t]
    \centering
    \includegraphics[width=0.9\linewidth]{Cap_3_figures/MOL_temperature/detuning_opt.pdf}
    \caption{Optical molasses temperature as a function of laser detuning. The horizontal axis represents the absolute value of the applied red detuning $|\delta|$, normalized to the natural linewidth $\Gamma/2\pi$. The data exhibits optimal sub-Doppler cooling in the region $|\delta|\approx 15-20$, reaching a minimum temperature of $13.9 \pm 0.4~\mu\mathrm{K}$. For detunings up to $|\delta|\approx 15$, the temperature follows the $1/|\delta|$ dependence expected for an ideal Sisyphus cooling. At larger detunings, the curve flattens and eventually begins to rise.} 
    \label{fig:detuning_molasses_temperature}
\end{figure}

The lowest average temperature was obtained for  a detuning of $|\delta|=17.2$ corresponding to $|\delta|=104.2$~MHz. Figure \ref{fig:11MHz_temp_fit} shows the fit of equation \eqref{eq:varince_expantion} for the optimal temperature configuration. The resulting temperature along the Cartesian axis are $T_x=(14.1\pm0.5)~\mu\mathrm{K}$ and $T_y=(13.5\pm0.7)~\mu\mathrm{K}$. Their weighted average yields the final temperature $T=(13.9\pm0.4)~\mu\mathrm{K}$.
\begin{figure}[h]
    \centering
    \includegraphics[width=0.8\linewidth]{Cap_3_figures/MOL_temperature/temperature_fit_10.8MHz.pdf}
    \caption{Optical molasses temperature fit for $|\delta|=17.2~\Gamma/2\pi$ with residuals. The large residuals of the last three data points from the linear fit are attributed to the cloud expansion, which makes its radial size non-negligible compared to the probe beam profile.}
    \label{fig:11MHz_temp_fit}
\end{figure}

As observed in Figure \ref{fig:11MHz_temp_fit}, the data points at long TOF diverge appreciably from the linear trend, consistently falling below the expected values. This behavior is consistently observed for all measurements. The hypothesis is that at sufficiently long TOF the cloud radius becomes non-negligible compared to the probe beam radius, violating the assumption of uniform illumination. This makes the outer region of the could lie in a region of lower intensity, which leads the Gaussian fit to systematically underestimate the true variance. Fortunately, this effect is expected to become less disturbing as the temperature is lowered.

%*###############################################################
\subsection{AOM amplitude optimization}
%*###############################################################
As a second optimization step, the influence of cooling beam power was investigated by varying the AOM drive amplitude and the duration of the amplitude ramp. The ramp is applied at the beginning of the molasses phase. The AOM amplitude expressed in arbitrary units (a.u.) controls the RF amplitude driving the AOM. Three values were tested: 385, 475 and 550~a.u., corresponding respectively to approximately 0.7, 0.85 and 1.0 times the AOM amplitude during the MOT phase. In addition, rap durations of  5~ms and 20~ms were tested for the transition of the AOM amplitude from the MOT value to the molasses value. The measurements were repeated for different values of the cooling detuning.
\begin{figure}[t]
    \centering
    \includegraphics[width=0.8\linewidth]{Cap_3_figures/MOL_temperature/aom_amplitude_comparison.pdf}
    \caption{Optical molasses temperature versus cooler detuning for different AOM amplitude and ramp durations. No clear optimal AOM amplitude is observed as the variations remain within the uncertainty. Temperatures corresponding to a ramp duration of 20~ms (dashed lines) are systematically lower than those obtained with 5~ms ramp (solid lines).}
    \label{fig:aom_amp_and_ramp_MOL_opt}
\end{figure}
The results of the optimization are shown in figure \ref{fig:aom_amp_and_ramp_MOL_opt}. Although small variations in the expected values of the temperature can be observed for different values of the AOM amplitude, these differences are smaller compared to the errors bars. Therefore, the measurements do not show statistically significant dependence of the molasses temperature on the AOM amplitude, and no clear optimum could be identified within the tested range. Hover lower amplitudes were not investigated to avoid operating the AOM too far from its optimal efficiency regime.

On the other hand, as shown in Figure \ref{fig:aom_amp_and_ramp_MOL_opt}, the temperatures obtained with a 20ms ramp are systematically lower that those obtained with a 5~ms ramp, yielding an average temperature reduction of approximately $0.5~\mu\mathrm{K}$. This improvement was relatively small, and was considered sufficient for the scope of the experiment, since further optimization of the order of few tens of microkelvin was non required, no other ramp times were tested. Consequently, a ramp duration of 20~ms was selected for the subsequent measurements.

It is worth noting that during the ODT loading characterization, this optimization was found to introduce additional fluctuations in the system. For this reason, the subsequent optimization of the compensation coil currents was carried out using an AOM amplitude of 385 a.u. corresponding to the lowest temperature measured at $|\delta|=17.2~\Gamma/2\pi$, with a ramp duration of 20~ms. However, during the later optimization of the ODT, these parameters were adjusting again. Nevertheless, the measurements presented in this section revealed that their influence on the molasses temperature remains limited to a few tens of microkelvin, and therefore does not significantly affect the final temperature achieved.


%*###############################################################
\subsection{Compensation field optimization}
%*###############################################################
The final stage of the molasses optimization procedure, consisted in adjusting the compensation magnetic field. The compensation coil currents are set to the molasses values 5~ms after the MOT magnetic field had been switched off. This 5~ms delay is introduced to ensure that the MOT magnetic field power supply had fully discharged, and the quadrupole field had sufficiently decayed.
The optimization was performed sequentially. First, the z-axis current was optimized while keeping the x- and y-axis currents fixed to the preliminary values \eqref{eq:shim_preliminary_values}. Then the x-axis current was optimized using the optimal z value, and finally the y-axis current was optimized using the optimal currents for both the z- and x-axis. The results are shown in figures \ref{fig:shimz_opt}, \ref{fig:shimx_opt} and \ref{fig:shimy_opt} for the z-, x-, and y-axis coils respectively.

The resulting optimized currents used in the subsequent ODT loading stage are:
\begin{equation}
    \begin{aligned}
        I_{\mathrm{x}} &= -15~\mathrm{mA} \\
        I_{\mathrm{y}} &= -185~\mathrm{mA} \\
        I_{\mathrm{z}} &= -230~\mathrm{mA}
    \end{aligned}
    \label{eq:mol_shim_optimal_values}
\end{equation}
The value of the z-axis compensation coil current deviated the most from its preliminary value, while the x-axis current underwent a smaller but still appreciably correction. In contrast, the y-axis coils current remained close to its preliminary values.

\begin{figure}[!h]
    \centering
    \includegraphics[width=0.75\linewidth]{Cap_3_figures/MOL_temperature/temperature_vs_shimz_z_fit.pdf}
    \caption{Optical molasses temperature versus. Shim-Z current. The parabolic fit retrieves a minimum temperature $T=(11.46\pm0.03)~\mu\mathrm{K}$ at a current $I_\text{z}=(-223.2\pm0.7)~\text{mA}$.}
    \label{fig:shimz_opt}
\end{figure}

\begin{figure}[!h]
    \centering
    \includegraphics[width=0.75\linewidth]{Cap_3_figures/MOL_temperature/temperature_vs_shimx_opt2_fit.pdf}
    \caption{Optical molasses temperature versus Shim-X current. The parabolic fit retrieves a minimum temperature $T=(9.2\pm0.1)~\mu\mathrm{K}$ at a current $I_\text{x}=(-17\pm2)~\text{mA}$.}
    \label{fig:shimx_opt}
\end{figure}
\begin{figure}[!h]
    \centering
    \includegraphics[width=0.75\linewidth]{Cap_3_figures/MOL_temperature/temperature_vs_shimy_opt2_fit.pdf}
    \caption{Optical molasses temperature versus Shim-Y current. The parabolic fit retrieves a minimum temperature of $T=(9.6\pm0.1)~\mu\mathrm{K}$ at a current $I_\text{y}=(-184\pm3)~\text{mA}$.}
    \label{fig:shimy_opt}
\end{figure}